% GENERATED FILE. Edit canonical lessons, then run npm run generate:books.
% canonical-source-hash: fnv1a64:503e185ebf615450
% canonical-lessons: MW-C188-phul, MW-C188-phal, MW-C188-nim, MW-C188-jhupri, MW-C188-chhapro

\chapter{Around the Courtyard}
\label{ch:mw-around-the-courtyard}
\hlchaptermodality{188}

\begin{chapteropening}
\textbf{By the end of this chapter:} \emph{I can name the flowers, trees and shelters around the house.}

Complete the last lesson of chapter 188: I can name the flowers, trees and shelters around the house.
\end{chapteropening}

\section[\texorpdfstring{\mw{फूल} (phūl)}{फूल (phūl)}]{\mw{फूल} (phūl) — a flower}

\label{lesson:MW-C188-phul}



\begin{quote}
Before the new one: say the Marwadi for a corner, then the Marwadi for a loft.
\end{quote}

\subsection*{You'll want to know: \mw{फूल}}
\textbf{\mw{फूल}} — \emph{phūl} — \textquotedblleft{}a flower\textquotedblright{}.

\begin{quote}
\mw{ये} \mw{फूल} \mw{लाल} \mw{है।}
\end{quote}

\begin{quote}
\emph{ye phūl lāl hai.}
\end{quote}

\begin{quote}
This flower is red.
\end{quote}

A masculine noun.

\begin{cousinweb}
\textbf{\mw{फूल}} (\emph{phūl}) is inherited from Sanskrit \emph{phulla}, \textquotedblleft{}blown (of a flower)\textquotedblright{}; CDIAL lists Old Marwari \emph{phūla} (CDIAL 9092).
\end{cousinweb}

\begin{grammarlens}[title={What you've built}]
One more word for this chapter.
\end{grammarlens}

\subsection*{Your turn}
\begin{itemize}
\raggedright
  \item \emph{Say it:} \emph{phūl}
  \item \emph{Say it:} \emph{phūl}, once more
  \item \emph{Say it:} \emph{ye phūl lāl hai.} — This flower is red.
  \item \emph{Recall:} say the Marwadi for a corner, then the Marwadi for a loft, then say \emph{phūl} again
\end{itemize}

\begin{tcolorbox}[breakable,colback=teal!4,colframe=teal!35!black,title={Before you move on}]
What does \mw{फूल} mean? (\textquotedblleft{}A flower\textquotedblright{}.) Say it once more.
\end{tcolorbox}
\section[\texorpdfstring{\mw{फळ} (phaḷ)}{फळ (phaḷ)}]{\mw{फळ} (phaḷ) — a fruit}

\label{lesson:MW-C188-phal}



\begin{quote}
Before the new one: say the Marwadi for a loft, then the Marwadi for a flower.
\end{quote}

\subsection*{You'll want to know: \mw{फळ}}
\textbf{\mw{फळ}} — \emph{phaḷ} — \textquotedblleft{}a fruit\textquotedblright{}.

\begin{quote}
\mw{ये} \mw{फळ} \mw{मीठो} \mw{है।}
\end{quote}

\begin{quote}
\emph{ye phaḷ mīṭho hai.}
\end{quote}

\begin{quote}
This fruit is sweet.
\end{quote}

A masculine noun.

\begin{cousinweb}
\textbf{\mw{फळ}} (\emph{phaḷ}) is inherited from Sanskrit \emph{phala}, \textquotedblleft{}fruit\textquotedblright{}; CDIAL lists Gujarati \emph{phaḷ} and Hindi \emph{phal} (CDIAL 9051). CDIAL gives no Marwari form in this entry.
\end{cousinweb}

\begin{grammarlens}[title={What you've built}]
One more word for this chapter.
\end{grammarlens}

\subsection*{Your turn}
\begin{itemize}
\raggedright
  \item \emph{Say it:} \emph{phaḷ}
  \item \emph{Say it:} \emph{phaḷ}, once more
  \item \emph{Say it:} \emph{ye phaḷ mīṭho hai.} — This fruit is sweet.
  \item \emph{Recall:} say the Marwadi for a loft, then the Marwadi for a flower, then say \emph{phaḷ} again
\end{itemize}

\begin{tcolorbox}[breakable,colback=teal!4,colframe=teal!35!black,title={Before you move on}]
What does \mw{फळ} mean? (\textquotedblleft{}A fruit\textquotedblright{}.) Say it once more.
\end{tcolorbox}
\section[\texorpdfstring{\mw{नीम} (nīm)}{नीम (nīm)}]{\mw{नीम} (nīm) — a neem tree}

\label{lesson:MW-C188-nim}



\begin{quote}
Before the new one: say the Marwadi for a flower, then the Marwadi for a fruit.
\end{quote}

\subsection*{You'll want to know: \mw{नीम}}
\textbf{\mw{नीम}} — \emph{nīm} — \textquotedblleft{}a neem tree\textquotedblright{}.

\begin{quote}
\mw{म्हारो} \mw{नीम} \mw{घणो} \mw{बड़ो} \mw{है।}
\end{quote}

\begin{quote}
\emph{mhāro nīm ghaṇo baṛo hai.}
\end{quote}

\begin{quote}
My neem tree is very big.
\end{quote}

A masculine noun.

\begin{cousinweb}
\textbf{\mw{नीम}} (\emph{nīm}) is inherited from Sanskrit \emph{nimba}, \textquotedblleft{}the neem tree\textquotedblright{}; CDIAL lists Hindi \emph{nīm} (CDIAL 7245). CDIAL gives no Marwari form in this entry.
\end{cousinweb}

\begin{grammarlens}[title={What you've built}]
One more word for this chapter.
\end{grammarlens}

\subsection*{Your turn}
\begin{itemize}
\raggedright
  \item \emph{Say it:} \emph{nīm}
  \item \emph{Say it:} \emph{nīm}, once more
  \item \emph{Say it:} \emph{mhāro nīm ghaṇo baṛo hai.} — My neem tree is very big.
  \item \emph{Recall:} say the Marwadi for a flower, then the Marwadi for a fruit, then say \emph{nīm} again
\end{itemize}

\begin{tcolorbox}[breakable,colback=teal!4,colframe=teal!35!black,title={Before you move on}]
What does \mw{नीम} mean? (\textquotedblleft{}A neem tree\textquotedblright{}.) Say it once more.
\end{tcolorbox}
\section[\texorpdfstring{\mw{झूंपड़ी} (jhū̃pṛī)}{झूंपड़ी (jhū̃pṛī)}]{\mw{झूंपड़ी} (jhū̃pṛī) — a hut}

\label{lesson:MW-C188-jhupri}



\begin{quote}
Before the new one: say the Marwadi for a fruit, then the Marwadi for a neem tree.
\end{quote}

\subsection*{You'll want to know: \mw{झूंपड़ी}}
\textbf{\mw{झूंपड़ी}} — \emph{jhū̃pṛī} — \textquotedblleft{}a hut\textquotedblright{}.

\begin{quote}
\mw{खेत} \mw{मांय} \mw{एक} \mw{छोटी} \mw{झूंपड़ी} \mw{है।}
\end{quote}

\begin{quote}
\emph{khet mā̃y ek chhoṭī jhū̃pṛī hai.}
\end{quote}

\begin{quote}
There's a small hut in the field.
\end{quote}

A feminine noun.

\begin{cousinweb}
\textbf{\mw{झूंपड़ी}} (\emph{jhū̃pṛī}) comes from a reconstructed \emph{*jhuppa}, \textquotedblleft{}a hut\textquotedblright{}; CDIAL lists Gujarati \emph{jhũpṛī} and Hindi \emph{jhopṛī} (CDIAL 5403). CDIAL gives no Marwari form in this entry.
\end{cousinweb}

\begin{grammarlens}[title={What you've built}]
One more word for this chapter.
\end{grammarlens}

\subsection*{Your turn}
\begin{itemize}
\raggedright
  \item \emph{Say it:} \emph{jhū̃pṛī}
  \item \emph{Say it:} \emph{jhū̃pṛī}, once more
  \item \emph{Say it:} \emph{khet mā̃y ek chhoṭī jhū̃pṛī hai.} — There's a small hut in the field.
  \item \emph{Recall:} say the Marwadi for a fruit, then the Marwadi for a neem tree, then say \emph{jhū̃pṛī} again
\end{itemize}

\begin{tcolorbox}[breakable,colback=teal!4,colframe=teal!35!black,title={Before you move on}]
What does \mw{झूंपड़ी} mean? (\textquotedblleft{}A hut\textquotedblright{}.) Say it once more.
\end{tcolorbox}
\section[\texorpdfstring{\mw{छपरो} (chhapro)}{छपरो (chhapro)}]{\mw{छपरो} (chhapro) — a thatched shelter}

\label{lesson:MW-C188-chhapro}



\begin{quote}
Before the new one: say the Marwadi for a neem tree, then the Marwadi for a hut.
\end{quote}

\subsection*{You'll want to know: \mw{छपरो}}
\textbf{\mw{छपरो}} — \emph{chhapro} — \textquotedblleft{}a thatched shelter\textquotedblright{}.

\begin{quote}
\mw{छपरो} \mw{जूनो} \mw{है।}
\end{quote}

\begin{quote}
\emph{chhapro jūno hai.}
\end{quote}

\begin{quote}
The thatched shelter is old.
\end{quote}

A masculine noun.

\begin{cousinweb}
\textbf{\mw{छपरो}} (\emph{chhapro}) is inherited from Sanskrit \emph{chattvara}, \textquotedblleft{}a house, a bower\textquotedblright{}; CDIAL lists Hindi \emph{chaprā}, \textquotedblleft{}a sloping thatch, a hut\textquotedblright{} (CDIAL 4976). CDIAL gives no Marwari form in this entry.
\end{cousinweb}

\begin{grammarlens}[title={What you've built}]
One more word for this chapter.
\end{grammarlens}

\subsection*{Your turn}
\begin{itemize}
\raggedright
  \item \emph{Say it:} \emph{chhapro}
  \item \emph{Say it:} \emph{chhapro}, once more
  \item \emph{Say it:} \emph{chhapro jūno hai.} — The thatched shelter is old.
  \item \emph{Recall:} say the Marwadi for a neem tree, then the Marwadi for a hut, then say \emph{chhapro} again
\end{itemize}

\begin{tcolorbox}[breakable,colback=teal!4,colframe=teal!35!black,title={Before you move on}]
What does \mw{छपरो} mean? (\textquotedblleft{}A thatched shelter\textquotedblright{}.) Say it once more.
\end{tcolorbox}
